\documentclass[10pt]{article}
\usepackage[a4paper,margin=0.67in]{geometry}
\usepackage{mathptmx}
\usepackage{microtype}
\usepackage{amsmath,amssymb,mathtools}
\usepackage{graphicx}
\usepackage{float}
\usepackage{placeins}
\usepackage{xcolor}
\usepackage{booktabs}
\usepackage{array}
\usepackage{enumitem}
\usepackage{caption}
\usepackage[hidelinks]{hyperref}
\hypersetup{pdftitle={From Point Scores to Certified Regional Probabilities},pdfauthor={Shuyin Xia; Guoyin Wang; Xinbo Gao}}
\usepackage[numbers,super,sort&compress]{natbib}
\setlength{\parindent}{0pt}
\setlength{\parskip}{2pt}
\setlist{nosep,leftmargin=*}
\captionsetup{font=small,labelfont=bf,justification=raggedright,singlelinecheck=false}
\definecolor{ink}{HTML}{1B2633}
\definecolor{blue}{HTML}{2F65D9}
\definecolor{teal}{HTML}{0B7A75}
\definecolor{orange}{HTML}{E76F24}
\definecolor{purple}{HTML}{6F3FCB}
\color{ink}
\newcommand{\R}{\mathbb{R}}
\newcommand{\E}{\mathbb{E}}
\newcommand{\Prob}{\mathbb{P}}
\newcommand{\ball}{\mathcal{B}}
\newcommand{\diam}{\operatorname{diam}}
\newcolumntype{L}[1]{>{\raggedright\arraybackslash}p{#1}}
\title{From Point Scores to Certified Regional Probabilities}
\author{Shuyin Xia$^{1,*}$, Guoyin Wang$^{2}$ and Xinbo Gao$^{3}$\\
$^{1}$School of Artificial Intelligence,\\
Chongqing University of Posts and Telecommunications, Chongqing 400065, China\\
$^{2}$Chongqing Normal University, Chongqing 401331, China\\
$^{3}$School of Electronic Engineering, Xidian University, Xi'an 710071, China\\
$^{*}$Correspondence: Shuyin Xia (xiasy@cqupt.edu.cn)}
\date{}
\begin{document}
\maketitle
\vspace{-0.4em}
\textbf{Summary.}
Scientific observations are usually recorded as points, whereas scientific decisions are often made over regions, groups or candidate sets. We introduce a regional probability object with explicit resolution; in this paper it is instantiated by a Granular Ball carrying an estimand, effective support, geometric extent, uncertainty, reuse conditions, refinement state and action status. This is not a new probability axiom: it changes the unit at which probability is computed and audited. Under the declared singleton/identity zero-radius branch, $r=0$ restores prediction, objective, calibration, decision, certificate state and fallback semantics of point Logistic. For $r>0$, a regional value is used only after support, homogeneity, boundary, calibration and transport certificates pass. A machine-readable head registry extends the same object contract to alternative links and uncertainty layers; each head must expose common certificates and a head-specific point-compatible fallback, so extensibility does not relax the failure boundary. Passing objects may be reused or counted as one downstream unit, and may be more stable within a declared perturbation envelope; failure triggers refinement, point fallback or abstention. The finite-sample result is conditional on independent validation and supplied structural envelopes. In the released audits, GB--Logistic--S was non-inferior to point Logistic on clean synthetic data and showed a conditional moderate-perturbation stability signal, while boundary crossing triggered fallback. STEAD reduced fixed-threshold inputs by 64.8\% but had no matched-recall call saving; LHC candidate compression carried assigned-event burden; Solar rejected regional action on its passing subset. These results define an auditable, extensible probability protocol with explicit limits rather than a universal accuracy claim.

\textbf{Introduction.}\;The usual statistical output of a probabilistic classifier is a point score attached to one observation. That output is useful, but it leaves the computational resolution implicit: it does not say which observations support the value, over what represented range it may be reused, how far its boundary extends, what uncertainty is induced by grouping, or what action follows when those conditions fail. Our question is therefore not whether metadata can be added to a point score, but whether a measurable region can itself be made the unit of probability computation while retaining a declared estimand and an exact route back to point inference. We do not alter the probability axioms. We expose a non-point probability object whose resolution, support, geometry, uncertainty and action contract are part of the statistical record. For a source measure $\mu_S$, the source regional estimand is $\eta_{A,S}=\E_{\mu_S}[Y\mid X\in A]$ (or $t_{A,S}=\E_{\mu_S}[\psi(Y)\mid X\in A]$); the deployment action estimand is $\theta_{A,D}=\int_A m_D(z)\,d\mu^{\rm deploy}_{A,D}(z)$, where $z=\phi(x)$ is the represented coordinate and $\mu^{\rm deploy}_{A,D}$ is a normalized target measure on $A$. They coincide only under an explicitly declared measure and conditional-law transport assumption.

We define the allocation cell first and the geometric representation second, but do not confuse geometric assignment with task quality. A region is a finite-resolution estimand, represented by a bounded envelope with a fixed tie rule and a refinement path. The record stores effective support, response summaries, geometry, uncertainty, refinement state and action status. The regional probability object is the statistical abstraction; a Granular Ball is the concrete, bounded and auditable geometric unit used here to instantiate it in a declared representation and norm (the primary implementation uses a fixed norm-ball). It is therefore more than a clustering label or a cache entry, while its semantics remain defined by the estimand and certificate rather than by geometry alone. Operationally, the object identity can be recorded as $\mathcal O_A=(\text{estimand},\text{allocation fingerprint},\text{support/geometry/uncertainty certificate},\text{head},\text{refinement state},\text{deployment/action contract},\text{fallback/failure state})$. Any change to the centre, radius, metric, tie rule, head or deployment measure creates a new object version and invalidates the old certificate and reuse record. A deterministic allocation induces the equivalence relation ``two represented states receive the same cell ID'', while a class may be carried by several balls and the number of balls is determined by local task difficulty rather than by the number of classes. Non-standard bounded envelopes (diagonal, low-rank, tree or other shapes) are compatibility and sensitivity extensions in the Supplementary Information, not separate semantic objects or an assumption that all shapes are interchangeable. The proof interface is allocator-agnostic: a GB-native allocator uses the declared ball split, centre, radius and refinement contract, whereas the released Solar, LHC and STEAD tables use KMeans- or MiniBatchKMeans-backed centres as regional allocation cells. Those KMeans-backed audits instantiate the interface; they do not establish that a KMeans optimizer is a GB-native optimizer or that either allocator dominates point gates.

\textbf{GB-native allocation audit.} To test the geometric identity directly, we also ran a bounded end-to-end audit with the released recursive Granular Ball allocator, using deterministic coordinate-median splits and a fail-closed support, boundary-margin and Brier-regret certificate. The GB-native-S head retained the point Logistic surface and added a support-shrunk local residual. Across 12 fixed synthetic seeds, 8.3\% of terminal balls (9.5\% of test queries) passed at zero label flipping; its Brier score was 0.19740 versus 0.19829 for point Logistic (mean difference $-0.00089$). With 20\% label flipping applied only to the structure/training role, only 1.3\% of balls (1.5\% of queries) passed and 98.5\% of queries returned to point Logistic; the S difference was $-0.00015$. This is an allocator and failure-boundary audit, not external evidence or a claim that native splitting universally dominates alternative allocators.

\begin{table}[H]
\centering
\small
\setlength{\tabcolsep}{3.5pt}
\renewcommand{\arraystretch}{1.16}
\begin{tabular}{@{}L{0.20\linewidth}L{0.47\linewidth}L{0.29\linewidth}@{}}
\toprule
Computational property & Existing interfaces (point, binning/cluster, hierarchical/local, conformal/set) & Regional probability object (GB implementation) \\
\midrule
Region as estimand & \textbf{Point:} point-only score. \textbf{Binning/cluster:} often an implicit cell summary. \textbf{Hierarchical/local:} a group or local target may be defined. \textbf{Conformal/set:} set coverage, usually not a regional mean. & Explicit source estimand $\eta_{A,S}$ and deployment action estimand $\theta_{A,D}$. \\
Effective support & \textbf{Point:} sample count is implicit. \textbf{Binning/cluster:} counts may be retained. \textbf{Hierarchical/local:} effective sample size or shrinkage. \textbf{Conformal/set:} calibration or coverage sample. & $n_{A,\mathrm{eff}}$, class balance and overlap are recorded. \\
Geometric boundary & \textbf{Point:} no region boundary. \textbf{Binning/cluster:} bin or cluster boundary. \textbf{Hierarchical/local:} parent or local boundary. \textbf{Conformal/set:} usually absent. & Granular Ball centre, radius, envelope, tie rule and boundary margin. \\
Regional uncertainty & \textbf{Point:} score or parameter uncertainty. \textbf{Binning/cluster:} optional within-cell variance. \textbf{Hierarchical/local:} posterior or partial-pooling uncertainty. \textbf{Conformal/set:} set-width or coverage uncertainty. & Heterogeneity, finite-sample, calibration, boundary and transport terms. \\
Refinement & \textbf{Point:} refit or rescore. \textbf{Binning/cluster:} split or recluster, often heuristic. \textbf{Hierarchical/local:} multiscale or local refit. \textbf{Conformal/set:} adaptive set size. & Split, re-certify and refit under a frozen contract. \\
Action contract & \textbf{Point:} external to the score. \textbf{Binning/cluster:} ad hoc threshold or queue. \textbf{Hierarchical/local:} model-dependent. \textbf{Conformal/set:} abstention may be available. & Deployment measure, budget and assigned-event semantics. \\
Point-compatible reduction & Native point output; other interfaces have special cases but generally no full semantic closure. & Under singleton/identity allocation, $r=0$ gives prediction, objective, calibration, decision, certificate and fallback closure; a rejected local gate gives point fallback, while zero local coefficients alone give point-surface and objective closure. \\
Certificate gate & Calibration, heuristic quality, posterior or coverage checks are usually separate interface checks. & Support, homogeneity, boundary, calibration, drift and transport gates are evaluated together. \\
Failure response & Outside-model, refit, shrink, ignore-cell or enlarge-set responses depend on the method. & Refine $\to$ point fallback $\to$ abstain. \\
Compression and stability & Compression and stability may occur, but are not intrinsic to the point, cell or set interface and require separate accounting. & Matched-risk/recall compression and assignment/reassignment perturbation audit. \\
\bottomrule
\end{tabular}
\caption{\textbf{Existing interfaces and the proposed regional probability object.} The table compares relevant statistical interfaces; it does not claim that the regional object is the only method with any individual property. The statistical abstraction is the regional probability object, while Granular Ball is the bounded, auditable geometric unit used to instantiate it. The head registry adds a common extension and failure interface rather than a model leaderboard; every registered head retains the point surface and inherits the same allocation, certificate, action and fallback contract.}
\label{tab:object-comparison}
\end{table}

Every experiment carries two distinct targets. Writing $z=\phi(x)$ for the represented coordinate, the sample-level target is the true conditional surface, $m_D(z)=\Prob_D(Y=1\mid \phi(X)=z)$, estimated by $\widehat q_A(z)$ and evaluated with row-level proper scores and calibration. The regional action target is instead $\theta_{A,D}=\int_A m_D(z)\,d\mu^{\rm deploy}_{A,D}(z)$, with estimator $\widehat\theta_{A,D}=\int_A\widehat q_A(z)\,d\widehat\mu^{\rm deploy}_{A,D}(z)$; it is used for a declared monitoring, review or budget decision under deployment measure $\mu^{\rm deploy}_{A,D}$. A regional action value is therefore not broadcast to every sample in $A$ as if it were a point probability: doing so incurs the within-region heterogeneity term and can misstate both Brier risk and calibration. The two targets share the frozen allocation record but have separate labels, measures, scores and promotion gates. The operating logic is consequently a single auditable chain rather than a collection of optional modules:
\[
\begin{gathered}
\text{regional probability object}\;\longrightarrow\;\text{certificate pass}\;\longrightarrow\;
\{\text{conditional reuse, compression, stability}\}\\
\longrightarrow\;\text{certificate failure}\;\longrightarrow\;
\{\text{refinement, point fallback, abstention}\}.
\end{gathered}
\]

The defining quality question is task-conditioned: does this region retain the information needed by the declared model and action at its current resolution, and can the task be represented by a simpler, sufficiently reliable regional model? For classification, this includes label support, class balance, purity or posterior width, within-region response heterogeneity and point--region Brier regret; for other tasks it uses the corresponding response and action loss. A different response mean alone is not a splitting rule. A region is refined when its current model is unreliable, unsupported or too complex for its certificate, and it is retained when a simpler model is adequate. Labels may enter both construction and quality: on the construction role they can guide class-aware seed quotas or boundary-directed candidates; on independent estimation and validation roles they determine response summaries and quality certificates. After selection, the allocator is frozen as a covariate-only map, so unseen queries require no outcome label and reuse remains within the declared contract. The released headline allocator uses label-free initialization; label-informed initialization is a candidate extension, not an established accuracy advantage. This separation makes task dependence explicit without turning the method into ordinary clustering or requiring one region per class.

The released allocators make a further geometric assumption that must be recorded rather than hidden: wherever Euclidean assignment is used, the coordinates are first standardized and the post-standardization feature weights are equal. Thus the current distance is an unweighted Euclidean distance in representation space, not a claim that all original physical variables are equally informative. A weighted candidate replaces it by $d_w^2(z,c)=\sum_j w_j(z_j-c_j)^2$ (or a regularized block/low-rank Mahalanobis form), with bounded weights and a fixed normalization. The weights are selected from a finite, predeclared menu on construction/calibration or validation roles using proper score, calibration, support and complexity gates; the selected metric, centres and tie rule are then frozen. This permits task-relevant geometry without making deployment label-dependent. Unless the grouped bootstrap interval passes the same non-inferiority gate as the regional probability, a learned metric remains a sensitivity or failure-boundary result rather than a precision claim.

For a label-informed candidate, the released protocol fixes the order rather than allowing labels to leak into deployment: (1) fit a label-free geometry-first parent partition that is shared across tasks; (2) use construction labels only to propose a finite set of class-aware seeds or task-conditioned distance terms; (3) estimate regional probabilities, Brier regret, calibration, posterior width and support on independent calibration/validation roles; (4) refine only parents with high out-of-fold residual, mixed boundary evidence and sufficient support; (5) freeze the resulting centres, cell IDs, distance and tie rule, so every query is assigned from covariates alone; and (6) audit both assignment-preserving and reassignment-aware perturbations with grouped bootstrap intervals before any candidate can replace the label-free headline allocator. The current release audits these components in separate candidate branches rather than claiming a jointly validated six-stage accuracy result. A candidate that fails these gates remains a failure-boundary diagnostic, irrespective of a favourable point estimate.

The central object is therefore a finite-resolution probability record with a point--region--point route. Projection and concentration arguments describe the regional estimand; independent refitting or a declared transport event returns to a point target only when its representation and geometric conditions are satisfied. Reuse is valid only under a frozen allocation contract. Failure is observable only through declared support, boundary, heterogeneity and transport triggers. Conditional stability must be tested with assignment-preserving or reassignment-aware perturbations; it is not assumed from a cached table. This makes the central proposition explicit: a region is provable under declared assumptions, reusable within a frozen allocation contract, auditable, failure-detectable through explicit triggers and conditionally stable.

The four consequences are implemented by a seven-stage regional-probability contract rather than seven independent accuracy claims. The contract starts with a frozen physical representation and a geometry-first parent allocation; it estimates an out-of-fold, support-conditioned parent--child residual; it integrates the resulting conditional surface over the declared region; it refines $K$ or a boundary only when local difficulty and support justify the split; it may borrow strength through an auditable region graph; it exposes a real region-first route with measured point calls and time; and, for multiple horizons, it may impose a joint monotone hazard calibration. Each extension records its estimate role, selection role, allocation fingerprint, uncertainty, support, failure reason and fallback. A new accuracy or cost statement is admitted only after independent calibration/validation, one outer test evaluation and a grouped-bootstrap gate (with ECE and coverage constraints); otherwise the extension remains a measured failure boundary in the Supplementary Information. This separation strengthens the interface itself: a regional probability object is not merely a coarser score, but a conditional computation with a declared estimand, a reusable state and an observable decision about when its local model must be rejected. The formal seven-module contract, certificate assumptions and conditional region-integral bound are given in SI S1. The theorem-to-experiment map labels each gate as measured, independently upper-bounded or model-assisted; the current Solar and simulation routes are partial/model-assisted gates rather than complete empirical validations of every theorem assumption. The SI separates conditional-surface error from within-region measure error and records a joint upper bound before any regional report is accepted. Sample-level predictions retain the continuous regional surface; regional action summaries integrate that surface against an explicitly declared target-population measure. A predeployment empirical measure is a proxy whose transfer requires overlap and conditional-law assumptions. In the new coupled weighting screen, both Solar horizons selected the unchanged point fallback; no incremental accuracy or deployment-measure identification claim follows. The staged model class used for interpretation is a hierarchy: a frozen physical point surface, a low-rank within-region residual, a deployment-measure estimator and a joint integration certificate precede any supervised split or route. The precision-directed follow-up screens retained the physical representation as the only confirmed point-anchor gain; continuous regional surfaces, support shrinkage, adaptive boundaries, low-support graph borrowing and conditional routing remain gated regional modules or measured failure boundaries.

For the precision-directed combination, these layers are executed in a strict seven-plus-one order. First freeze the validated physical point representation and parent allocation; then keep the sample-level surface $\widehat q_A(z)$ separate from the region-level action estimand $\theta_{A,D}=\int_A m_D(z)\,d\mu_{A,D}^{\rm deploy}(z)$ and its estimator $\widehat\theta_{A,D}=\int_A\widehat q_A(z)\,d\widehat\mu_{A,D}^{\rm deploy}(z)$; estimate the target deployment measure on an outer-time development/calibration role; fit the physical baseline plus a low-rank within-region residual; learn its support- and posterior-width-conditioned prior; refine only high-heterogeneity, adequately supported boundaries; and gate the regional output using total risk and the joint surface--measure uncertainty radius. Only after these seven checks pass is the eighth operation enabled: a risk-gated region-first route with measured point calls, wall time and downstream cost. Graph borrowing and joint hazard calibration are optional auxiliary operators within this sequence, not substitutes for the estimand, measure or uncertainty gates. The current release does not pass the regional-integral gate, so the final route is reported as an audit rather than an accuracy or efficiency claim.

An orthogonal region-level estimator remains an unvalidated follow-up: a cross-fitted, density-ratio-weighted residual correction for $\theta_{A,D}$ with overlap, effective-support and ratio-clipping diagnostics. It is specified to use region-action scores and grouped influence-function bootstrap, while $\widehat q_A(x)$ remains evaluated with sample-level Brier and ECE. This separates surface and deployment-measure nuisance terms; it does not assume that an unlabeled selector window identifies the future deployment law, and no result from this estimator is promoted here.

A calibration-only surface--action gate on the frozen $K=32$ allocation activated 47.7\%/39.0\% of Solar rows at 24/48 h and changed sample-level Brier by only $+0.000006/+0.000006$, with intervals crossing zero. Its separate regional-action scores are not comparable with sample-level Brier, so the gate is retained as a fallback and estimand-separation audit rather than a precision claim.

For Solar, the deployment label is defined at the HARPNUM--12-h-anchor--horizon level: an observation is positive when the first eligible M1.0-or-higher event occurs strictly after the anchor and within the declared 24/48-h window. The frozen protocol specifies that any outer-time analysis must estimate the deployment measure and a cross-fitted doubly robust regional integral only from earlier development/calibration blocks, then evaluate once on a genuinely unseen time block. The released 2019--2024 endpoint is currently protocol-only, so no outer-time DR accuracy is promoted in this release.

Resolution and initialization screens changed row-wise Brier in historical replays, but the point head remained better and neither screen identified a regional or action gain. They are therefore SI-only stability diagnostics. The staged surface--measure audit freezes $K=32,n_{\rm init}=1$, while $K=64$ remains the clean benchmark; a resolution change is a new allocation version and is admissible only after the same estimand, deployment-measure and outer-time gates are rerun.
\begin{figure}[t]
\centering
\includegraphics[width=0.96\linewidth]{figures/fig1_hero.pdf}
\caption{\textbf{Non-point computation as a measurable hierarchy.} A region stores probability-relevant mass and response summaries together with its geometric scale. The hierarchy is refined only where evidence changes. The certificate keeps selection, geometry, representation, dependence and implementation costs separate. The resulting regional table is a reusable interface for routing, calibration and review. The coarse-to-fine motivation is a design analogy to global-before-local perception, not a claim about a literal neural implementation.}
\end{figure}
\FloatBarrier

\textbf{Three contributions and four tests.}\;The first contribution is a new probability object: a region is an estimand-bearing record with effective support, geometry, uncertainty, refinement state and action semantics, rather than a cluster label, bin or cached point score. The second is a unified point--region--point protocol: under the declared singleton/identity branch, point Logistic is the strict zero-radius boundary, positive-radius reuse is certificate-gated, and certificate failure has a defined refinement, point-fallback or abstention route. A finite head registry enforces this protocol across computational heads, so closure, gates and fallback are inherited rather than redefined per model. The third is a set of conditional scientific capabilities: when support, homogeneity, boundary and transport conditions pass, several observations may share one regional probability and action state, potentially reducing downstream objects or calls and limiting degradation under a declared perturbation budget; when they fail, the same object exposes the cost and the failure boundary. The Supplementary Information proves a finite-menu certificate-sufficiency theorem and reports both a synthetic audit and a held-out Solar temporal test role. In that Solar replay, 7/32 cells pass at 24 h (30.01\% of test rows; 69.99\% fallback) and 4/32 at 48 h (17.57\%; 82.43\% fallback); the passing-subset regional action risk is not lower than its matched point risk, so the present result is a rejection/fallback boundary rather than a precision promotion. The four testable consequences are object reuse, matched-risk or matched-recall resource accounting, perturbation stability within the certificate envelope and explicit failure exit; failure and cost are part of the released object.



\textbf{Conditional resource and robustness.}\;We report certificate passage, fallback, passed-subset risk, the matched point baseline and a declared resource measure. The multi-seed simulation shows that object reuse is conditional: heterogeneous cells provide a small descriptive Brier change, whereas homogeneous, boundary and contaminated cases mostly fall back. A dedicated 20-seed GB--Logistic--S audit tests the nested soft head directly. On clean data its fail-closed route is below point Logistic by $-0.003084$ in Brier; at a $0.25$ assignment-margin perturbation it remains below point by $-0.003129$, with clean-to-perturbed Brier increases $+0.000014$ for the route and $+0.000059$ for point. On the same certified rows, cached S has zero Brier degradation versus $+0.000072$ for point; the paired route-minus-point slope is $-0.000600$ with descriptive seed-bootstrap interval $[-0.000957,-0.000228]$. These intervals are model-assisted stress evidence, not theorem coverage. At half-margin the corresponding matched point degradation is $+0.000543$. Directed boundary crossing changes 38.34\% of assignments and sends 100\% of rows to point fallback; large perturbation changes 47.75\% and leaves 99.90\% on fallback. Thus S showed no clean-data degradation in the descriptive synthetic screen, a conditional moderate-perturbation stability signal, and fail-closed behavior when its allocation contract was violated. The historical Solar action gate still rejects regional use on its passing subset, while STEAD and LHC expose matched-recall and assigned-event costs. These negative results define the method's scope rather than being converted into universal robustness or efficiency claims.


\textbf{Canonical head and extensibility.}\;Table~\ref{tab:object-comparison} summarizes the object-level distinction; the registry instantiates the same object contract with alternative probability heads without changing allocation, certificates, action or fallback semantics. GB--Logistic--S is the canonical strongly coupled computational head: it preserves the point surface and adds only certificate-gated soft local corrections. GB--Logistic--H remains the hard support, boundary, audit and fallback control. Point Logistic is the strict singleton/identity zero-radius baseline; setting local terms to zero gives point-surface and objective closure, while full decision, certificate and fallback closure requires the declared zero-radius branch or an explicitly rejected local gate. The registry in \texttt{evidence/head\_registry/} defines four registered point-compatible extensions—GB--Logistic--PP, GB--Beta--Binomial--S, GB--cloglog/GB--Hazard--S and GB--Dirichlet--S—using the same allocation fingerprints, support, heterogeneity, boundary, calibration, transport and action fields. Each declares its surface-preserving or corresponding point-hazard/softmax closure and a fail-closed point route. The registry is also an invalidation ledger: changing the allocation fingerprint, deployment measure or head-specific certificate invalidates the associated reuse record. The main text admits at most one extension after clean-data non-inferiority, independent certified-subset benefit, no false-pass increase, correct boundary triggering and quantified assigned-work cost. None of the four extensions passes all gates in this release; they remain extensibility protocols and failure boundaries rather than additional accuracy claims.

\textbf{Results.}\;The main text contains the regional probability object, its point--region--point theory, the four consequences and the seven-stage contract; the STEAD downstream-call test; the solar temporal model/evaluation audit plus a frozen external protocol; the LHC candidate-region boundary audit; and a unified failure boundary. EEG and digits remain supporting diagnostics in the Supplementary Information, and weather remains protocol-only. The experiments report proper scores, support, geometry, selected objects, assigned observations, downstream calls, resource scope and failure states under their declared denominators. No result is promoted to a universal accuracy, cost, energy or robustness claim.

\textbf{Certificate scope and experiment map.}\;The singleton/identity zero-radius theorem is an exact executable identity check for prediction, objective, Dirac target, calibration and action. The finite-menu sufficiency theorem is a conditional result for a frozen allocation and independently validated gates; the current empirical gates are explicitly v1 partial/model-assisted implementations, with transport and deployment-measure terms treated as assumptions or failure flags. The four domains therefore have distinct roles: the Solar temporal table tests rejection and point fallback; LHC tests candidate-object compression against assigned-event burden; STEAD tests input reduction at fixed recall and reports the absence of a matched-recall call saving; EEG is a supporting audit of regional reuse and coverage. The grouped audit makes these boundaries explicit: across 236 station-hash resamples, STEAD regional Brier exceeded point Brier by $+0.01894$ [$+0.00999,+0.02970$] and the 80\% recall queue assigned $+34$ traces [$+8,+74$]; the 19-subject EEG audit reports granular-region subject Brier $0.24883$ [$0.23816,0.25898$] but empirical interval coverage $0.516$; Solar HARPNUM exports support action counts only, not a group-risk interval. These are grouped audits and failure boundaries, not universal accuracy or efficiency claims. A predeclared synthetic coverage grid crosses effective support, $K$, homogeneity, boundary $b/\pi$, response and deployment-measure drift, and finite menu size $M$: within declared envelopes, interval coverage was 1.000 with no selective-risk violation in the retained cells; when the true envelopes were 2.5 times the declarations, coverage fell to 0.955 and a noisy-diagnostic false-pass stratum reached 10.6\%. This is an empirical proof-contract audit, not distribution-free or external-data coverage. The released grid is the constant-head subcase ($u_p=u_\mu=0$ and zero joint surface--measure remainder); non-zero surface/measure remainders remain theorem-level or model-assisted. Synthetic data also test conditional heterogeneity and contamination patterns. This ledger separates theorem guarantees from empirical claims and records which certificate version is being evaluated.


\textbf{Nested GB--Logistic heads.}\;The fixed-allocation heads form an explicit nested family. With $a(z)$ denoting the frozen cell and $r_j(z)$ a support-shrunk local coordinate,
\[
\eta_S(z)=\eta_{\rm point}(z)+g_{a(z)}\gamma_{a(z)}r_{a(z)}(z),\qquad g_j\in\{0,1\},
\]
\[
\text{point Logistic}\;\subset\;\text{GB--Logistic--S},\qquad
 g_j\gamma_j=0\ \forall j\ \Longrightarrow\ \widehat p_S(z)=\widehat p_{\rm point}(z).
\]
\noindent Thus zero local coefficients give point-surface/objective equality, while every local gate rejected gives the declared point fallback; full certificate and action equality is supplied by the singleton/identity zero-radius branch. GB--Logistic--H retains the hard cell for support, audit and fallback; S adds a gated soft residual only on calibration-approved cells. The frozen Solar head probe found a conditional S--H Brier difference of $-0.000655$/$-0.000606$ at 24/48 h, while S versus the matched point head crossed zero; the result is therefore a nested-interface and failure-boundary result, not a general precision claim. A centre or radius change creates a new allocation version and requires new certificates.

\textbf{A region as a probability object.}\;Let $\phi(X)\in\R^d$ be a representation frozen on a structure sample and let $\Pi_h=\{A_1,\ldots,A_K\}$ be a finite measurable partition. For a bounded score $\psi(Y)$, define $\mu_j=\Prob(\phi(X)\in A_j)$ and $t_j=\E[\psi(Y)\mid \phi(X)\in A_j]$. The conditional expectation $t_{\Pi_h}(X)=\sum_jt_j1\{\phi(X)\in A_j\}$ is an orthogonal projection. If $\Pi_{h'}$ refines $\Pi_h$, then
\[
\E[(\psi(Y)-t_{\Pi_h})^2]-\E[(\psi(Y)-t_{\Pi_{h'}})^2]=\E[(t_{\Pi_{h'}}-t_{\Pi_h})^2]\geq0.
\]
Thus refinement cannot increase population approximation error, although finite-sample variance can increase. For nested partitions generating $\sigma(\phi(X))$, the limit in $L_1$ is $\E[\psi(Y)\mid\phi(X)]$. It equals $\E[\psi(Y)\mid X]$ only when the representation preserves this conditional mean (SI S2). A finite-radius region is therefore a target at a declared resolution, and the point target is a limit rather than the only admissible object. The same chain accounts for representation loss, shape entropy, honest refitting, geometric transport, group dependence and optimization remainders; this article tests what that object enables when exposed as a reusable table.

Under iid groups, a positive denominator, a nondegenerate variance and the stated moment conditions, independent refitting gives asymptotic normality for the regional target. Point centring additionally requires undersmoothing. The formal result conditions on a realized cell count $N_A=k_n\to\infty$ (or assumes this count condition holds with probability tending to one); it does not assert unconditional fixed-$x$ normality. For the group-balanced ratio target, with $B$ independent refit groups and the group influence-score variance from SI S9, the studentized estimator obeys
\[
\frac{\sqrt{B}(\hat t_A-t_A)}{\hat\sigma_A}\Rightarrow N(0,1),
\]
while a point-target interval additionally expands by a valid geometric envelope. This establishes the inferential distinction that is easy to lose in a point-based implementation: a nominal region interval can be exact for $t_A$ and still undercover $t(x)$ when $\rho_A^\alpha$ is not controlled.

\textbf{Solar activity-region audit.}\;The solar evidence has two explicitly separated time layers. The model and evaluation state table is confined to SWAN--SF data from 2010--2018, with the test period in 2014--2018. The requested 2019--2024 JSOC/CGEM HMI/SHARP external endpoint is protocol-only in this release: the external input, independent GOES join and sampled FITS receipts are absent, and the bundled 20-row fixture is schema validation only. No temporal external metric, provider-field parity or raw-FITS claim is used for refitting, threshold selection or action claims. The release nevertheless includes a small fail-closed transport-audit interface for sampled FITS receipts, provider-versus-raw feature parity and independent GOES--HEK association; with no supplied raw manifest it reports pending data and produces no predictive metric. The frozen protocol retains the anchor estimand, two-hour lag transform, model, $\ell_1,K=64$ geometry, thresholds, tie rules and refinement contract, and requires positive M1.0+ labels, complete follow-up, explicit HARPNUM/NOAA\_AR association, quality and missingness fields and input hashes. The primary target is M1.0-or-higher (M+) flares in the next 24 h or 48 h; X+ is secondary and sparse. The state table contains 27,745 twelve-hour decision anchors, with HARPNUM-disjoint chronological roles; state-label counts are not independent flare-event counts.

The solar task is defined as allocation under a finite satellite, telescope and analyst budget. The raw object is an ordered activity-region time window; a Granular Ball $A$ is a measurable state region whose members have similar flux, shear, helicity, gradient and short-term evolution features. We estimate
\begin{equation}
p_{A,S}^{(h)}=\Pr_{\mu_S}\{Y_{M+}^{(h)}=1\mid X\in A\},\qquad n_{A,\mathrm{eff}}=N_{\mathrm{eff}}(A).
\end{equation}
Here $n_{A,\mathrm{eff}}$ is the number of valid, de-duplicated HARPNUM--time-window observations supporting the region. The source conditional mean $p_{A,S}^{(h)}$ is not by itself the deployment action target $\theta_{A,D}^{(h)}$. Its geometry is the centre, declared norm/shape, standardized scale, an optional principal direction, and a 95\% support boundary in the declared SHARP state coordinates; this is a magnetic-state geometry, not a claim about a solar-surface boundary. The action output is a queue of physical HARPNUM states for continued monitoring, rather than a requirement to inspect every timestamp.

The two quantities are scored separately. The sample-level query $q_A^{(h)}(x)$ is evaluated by Brier score, log loss and ECE. The action-level target is the same conditional surface integrated against a normalized, declared deployment measure on the region,
\begin{equation}
\theta_{A,D}^{(h)}=\int_A q_A^{(h)}(x)\,d\mu_{A,D}^{\rm deploy}(x),\qquad
R_h(b)=\frac{\sum_gY_{g,h}a_g(b)}{\sum_gY_{g,h}},\qquad
C_h(r)=\min\left\{N_{\rm win}:R_h(b)\ge r\right\},
\end{equation}
where $g=(\mathrm{HARPNUM},\mathrm{UTC\ day})$ is the deduplicated action unit, $a_g(b)$ selects the first $b$ fraction of units, $R_h(b)$ is positive M+ HARPNUM--day recall and $C_h(r)$ is the minimum selected-window burden at recall $r$. False alarms per active-region day and cost per captured positive are reported alongside these curves. A regional integral is never broadcast back to individual states; it is enabled only when its deployment measure, support, heterogeneity and uncertainty certificates pass. The anchor table has no independent GOES event identifier, so positive HARPNUM--day is an event-recall proxy; the bootstrap resamples HARPNUM/time blocks, and state counts are not interpreted as independent flare events. M+ is the primary endpoint; X+ is a descriptive sparse stratum because only nine positive HARPNUMs occur at each horizon.

The historical point-representation repair changed Brier by $-0.000334$ (interval $[-0.000633,-0.000063]$) at 24 h and $-0.000468$ (interval $[-0.000801,-0.000139]$) at 48 h. It is a point-anchor result, not a regional gain; the matched regional action route therefore remains fail-closed. The Solar action target uses deduplicated HARPNUM--UTC-day units and a declared monitoring measure. In the K32 conditional gate, 7/32 cells pass at 24 h (30.01\% of rows) and 4/32 at 48 h (17.57\%); the passed-subset regional risk is higher than its matched point risk, so this is an explicit Solar failure boundary. The strict K64 outer-time route likewise selects no safe cells and makes 100\% point calls. K32/K64, splits and estimands are fixed in the Solar run registry; historical resolution, initialization and cached-noise probes remain SI-only diagnostics.

The Solar evidence thus separates three objects: a clean K64 row-probability benchmark, a conditional K32 action certificate and a strict outer-time fallback route. A regional value is never broadcast to every timestamp; it is integrated against the declared deployment measure only after support, heterogeneity, boundary and transport gates pass. The released 2019--2024 raw-FITS/GOES endpoint remains protocol-only; the bundled transport interface is a pending-data handoff rather than an external result, and $N_{\rm HR}^{\rm proxy}$ is a derived SHARP-analysis proxy rather than a satellite-read count.

\begin{figure}[t]
\centering
\includegraphics[width=0.92\linewidth]{figures/solar_flagship_audit.pdf}
\caption{\textbf{Solar active-region state cost and conditional stress audit.} Panels a,b show distinct HARPNUM and 12-h state-window counts at matched positive-state M+ recall for 24- and 48-h horizons; SHARP-derived high-resolution reads are proxies and raw HMI reads were not measured. Panels c,d show Brier-score changes under sensor spikes and random-label contamination with clean $K=64$ geometry and cached assignments, using one seed. The regional object exposes support and geometry, but the point logistic model retains better clean proper scores and ranking in this split; the stress panels are conditional diagnostics, not general robustness evidence. The 2019--2024 provider-product endpoint is protocol-only; full raw-FITS/GOES reconstruction remains pending. X+ remains secondary because its test positives are sparse.}
\end{figure}
\FloatBarrier

\textbf{LHC candidate-region audit.}\;In LHC Olympics BlackBox1, the declared question is whether $10^6$ events can be represented by a small phase-space menu without hiding assigned-event burden. We freeze the three-dimensional state $z=(\log(1+m_{j,1}),\log(1+p_{T,j_1}),\log(1+m_{\rm all}))$, a $K=64$ menu, sideband rule and development/audit split before opening the audit stream. At target recalls 50/70/80/90/95\%, point Isolation Forest reviews 74,951/108,955/147,151/194,197/225,360 events, whereas regional prefixes select 28/53/53/53/59 candidate objects but assign 137,656/242,043/242,043/242,043/277,107 events. At common background-FPR caps 0.025/0.05/0.10, point recalls are 0.008/0.042/0.221 from 7,495/14,997/30,031 selected events, while regional prefixes select 1/3/5 objects, assign 6,243/13,669/26,632 events and recall 0/0.084/0.087. Candidate-object compression can trade against assigned-event burden and discrete recall jumps; there is no uniform matched-operating-point saving. A lightweight cross-split sideband-transfer probe over the 64 frozen cells gives a signal-subtracted closure-ratio median of 0.9934 (5th--95th percentile 0.9044--1.0993; support-weighted ratio 0.9909), and the feature/assignment/stress audit takes 2.9--27.2~s across its declared mechanisms. A development-only support-screen replay at target recall 0.90 changes the frozen prefix from 53/242,043 to 49/235,402 selected regions/assigned events at support $\geq 5{,}000$ (achieved recall/FPR 0.9011/0.7846); this is a descriptive post-freeze sensitivity row. The separate P4 aggregate audit propagates a counting-only independent-Poisson delta approximation and decomposes assigned events by fixed cell strata, but does not add prospective validation. These are descriptive processing and background-model checks: the original probe is not background-only and omits transfer-factor uncertainty, while detector/systematic effects, downstream fit and analyst time, look-elsewhere correction and discovery-level global $p$-value remain unmeasured. The full closure and physics-fit extensions remain in the Supplementary Information; no row is called a discovery.\cite{kasieczka2021,bump_hunter}

\begin{figure}[t]
\centering
\includegraphics[width=0.96\linewidth]{figures/lhco_region_audit.pdf}
\caption{\textbf{Candidate-region and assigned-event costs in the LHC Olympics benchmark.} The official event-level rerun fixes the three-dimensional representation, sideband rule and $K=64$ cells. Target and achieved recall are distinguished because whole-region selection has discrete jumps. Candidate counts and assigned-event counts are different units: fewer selected regional objects coexist with higher assigned-event FPR. These curves do not time downstream physics fits, establish a discovery-level global $p$-value or demonstrate a new particle. The separate six-mechanism stress audit uses 240 recorded rows; smearing and pileup are derived-feature surrogates, while label and duplicate mechanisms follow their declared frozen protocols.}
\end{figure}
\FloatBarrier

\textbf{STEAD stability and object audit.}\;We checksum-verified a 12,000 three-component waveform-window STEAD subsample and held out 2,381 PhaseNet inputs by station hash. Twelve frozen waveform descriptors feed a point logistic gate, a score-bin gate, a depth-six tree, a $K=64$ regional allocation table implemented by quality-weighted MiniBatchKMeans and a diagonal-Gaussian mixture comparator. At earthquake-window recall 0.80, the five queues select 815, 835, 808, 839 and 870 waveform windows, respectively; at 0.90 they select 936, 952, 923, 974 and 982. The same PhaseNet picker receives all 2,381 held-out waveform windows or 839 high-risk waveform windows in 23 regional cells, giving a 64.8\% input reduction at recall 0.80. The explicit matched-recall audit gives 815 point traces versus 839 assigned regional traces at exactly 0.80 recall; at nominal recalls 0.70/0.90/0.95, the regional prefixes assign 736/974/1,068 traces versus 711/936/999 point traces, with discrete regional recall 0.7168/0.9050/0.9525. At matched observed FPR 0.20, the regional prefix assigns 1,242 waveform windows at recall 0.967, while point, score-bin, tree and GMM queues assign 1,274, 1,255, 1,280 and 1,257 waveform windows at recalls 0.990, 0.990, 0.996 and 0.974. Thus the 64.8\% reduction is a fixed-threshold coverage--call trade-off, not a matched-recall call advantage; selected regional cells do not imply fewer assigned waveform calls. Clean Brier/log loss/ECE are 0.0330/0.1343/0.0256 (point), 0.0340/0.1318/0.0228 (score-bin), 0.0265/0.1059/0.0122 (tree), 0.0524/0.1930/0.0345 (regional allocation) and 0.0556/0.2021/0.0389 (GMM). A real CPU PhaseNet original-v2 call measured 28.10 s for all 12,000 waveform windows, 5.49 s for the 2,381-waveform-window holdout and 2.05 s for the 839-waveform-window regional queue (23 whole cells; batch size 256). The same receipt measured 0.229 s raw HDF5 loading, 7.78 s frozen feature extraction, 0.013 s model loading and 0.031 s queue selection; peak RSS was 1.94 GiB. The 15-W number is a wall-time energy proxy, not instrumented power, and these endpoints do not establish PhaseNet accuracy, operational warning utility or a universal matched-recall saving. A separate real-domain GB-native audit on the same checksum-verified 12,000-window subset passed 1/64 terminal balls, routed 3.61\% of holdout queries regionally and sent 96.39\% to point fallback; GB-native-S Brier was 0.028312 versus 0.028269 for point Logistic. This is identity and failure-boundary evidence, not a general accuracy or resource gain.

The perturbation tables are mechanism-specific. Cached assignments diagnose the frozen-allocation contract; feature-driven reassignment, support loss, boundary crossing or uncontrolled transport remainder invalidate reuse and trigger refinement or an unsupported status. The raw-waveform stress is a single-seed subsample audit, not a universal seismic robustness result.

\begin{figure}[t]
\centering
\includegraphics[width=0.98\linewidth]{figures/stead_raw_audit.pdf}
\caption{\textbf{Raw-waveform STEAD subsample robustness audit.} Station-hash holdout on 12,000 checksum-verified waveform windows compares selected point waveform windows with supported $K=64$ waveform cells at matched recall and matched observed background-FPR caps; the stress panel shows the channel-dropout mechanism after feature re-extraction; the full six-mechanism table is in the Supplementary Information. Regions reduce candidate objects but can assign more waveform windows and have weaker clean proper scores. The perturbation curves are single-seed, mechanism-specific diagnostics on a Zenodo subsample, not a full-release or universal seismic-robustness claim.}
\end{figure}
\FloatBarrier






\textbf{Discussion.}\;The evidence supports a finite-resolution probability object rather than a point score with an audit label attached after the fact. Its record binds an estimand to effective support, Granular Ball geometry, uncertainty, refinement state and action contract. The central chain is \emph{object $\to$ certificate $\to$ conditional reuse/compression/stability $\to$ failure exit}: under the stated sampling, selection, representation, transport and perturbation assumptions, the certificate theorem gives conditional regional guarantees; within a frozen allocation contract, the probability and computation state can be reused; when support, homogeneity, boundary or drift fails, the object carries the reason for refinement, point fallback or abstention. The GB--Logistic--S stress audit supplies a model-assisted conditional stability signal, while Solar, STEAD and LHC show where certificates reject regional reuse or where apparent object compression carries a cost.

The statistical abstraction is a regional probability object and the implementation contribution is a Granular Ball as its bounded, auditable geometric unit. The registry tests extensibility at the protocol level: a new link or variance layer qualifies only if it preserves the point surface, declares its closure, inherits the certificate/action/fallback fields and passes independent promotion gates; adding a link function alone is not a new object. This division preserves generality without erasing the method identity: the Granular Ball supplies centre, radius, boundary margin, allocation fingerprint and refinement path, while the certificate supplies the conditions under which that geometry may carry a probability and an action state. The exact singleton/identity zero-radius closure makes the compatibility claim stronger than a loose metadata analogy: $r=0$ under singleton/identity allocation recovers point prediction, objective, calibration, decision, certificate status and fallback semantics. A rejected local gate gives point fallback, whereas zero local coefficients alone establish point-surface and objective equality. The present evidence establishes neither a particle discovery, a solar-flare mechanism, clinical validity, operational earthquake warning, universal energy saving nor assumption-free pointwise coverage. It establishes a falsifiable way to change probability resolution while preserving a strict point baseline, exposing conditional resource and perturbation consequences, and making failure an explicit output.

\textbf{Statistics.} Solar uses chronological and HARPNUM-disjoint roles and reports state-window, HARPNUM and AR-day denominators separately. LHC uses a development/audit split and descriptive stress rows; candidate objects, assigned events and downstream fits are not pooled. STEAD uses a station-hash holdout and one frozen PhaseNet weight; its whole-cell station bootstrap gives regional-minus-point Brier 0.0189 [0.0100, 0.0297], log-loss 0.0576 [0.0295, 0.0895] and assigned-trace difference 34 [8, 74] at the 0.80 endpoint. These are descriptive frozen-prediction intervals. EEG resamples query subjects, solar resamples perturbation seeds within fixed blocks and LHC reports ranges because repeat identifiers are non-unique.

\textbf{Methods.} The reference implementation uses disjoint structure, estimation, validation, refit, query and reference streams where required. Each regional record stores allocation, support, geometry, uncertainty, refinement, action status and failure reason. Empirical support-relative position and distance to the actual frozen allocation boundary are recorded separately; the latter is compared with a perturbation budget in the same declared metric before reuse. Relative distance changes only quality auditing and routing unless a new allocation, probability table and version are explicitly fitted and frozen. Assignment-preserving perturbations retain the membership map and test cached-state validity; reassignment-aware perturbations recompute membership and test boundary, support and transport. Any invalidation requires recomputation or an unsupported status. The Supplementary Information contains the assumptions, proofs, uncertainty intervals and source-table paths. A fixed-assignment audit additionally records relative support scale, exact all-centre allocation margin, perturbation budget and calibration residual risk; these quantities route or reject reuse and do not redefine the partition unless a new probability table is refitted and frozen.

The resolution and initialization screens were held separate from the precision target: increasing $K$ or $n_{\rm init}$ can change allocation stability and support, but it cannot repair a mismatched estimand, within-ball aggregation or deployment measure. The staged surface--measure audit therefore freezes $K=32$ and $n_{\rm init}=1$, while $K=64$ remains the clean benchmark; any later change in resolution must be re-estimated and re-gated on its declared roles.

\textbf{Data availability.}\;The release provides promoted summary tables, split indices, figure-source data, query manifests, returned counts and SHA-256 checksums; the promoted SWAN summaries are linked to the cited public record; the full state table is not bundled and a complete rerun requires the full archive. Large public sources are not duplicated; download scripts identify the official records and verify their hashes. The solar model/evaluation table is confined to SWAN--SF 2010--2018; the full derived anchor table is not bundled; a complete rerun requires the cited SWAN archive, while the included script supports deterministic subset verification. The requested 2019--2024 HMI/SHARP external endpoint remains protocol-only in this release: an auxiliary event-only GOES/HEK provenance probe is included, but the raw-FITS-linked event join, provider-field parity, sampled FITS receipts and predictive external metrics are not present or claimed. The digits audit uses \texttt{sklearn.datasets.load\_digits} (1,797 samples, 64 features and 10 classes). The full raw-FITS/GOES endpoint, full STEAD parent archive and large LHC HDF5 files are not claimed as bundled reruns. The release tag, license and canonical PDF hashes are fixed in the package manifest; the corresponding author must add the permanent repository URL and DOI after deposit before public-release claims are made.

\textbf{Code availability.}\;The tagged code release contains the analysis scripts, frozen seeds, baseline implementations, figure-generation scripts, environment lockfiles, derived tables and immutable SHA-256 manifests. The reproducibility runner resolves public inputs, checks their hashes and regenerates the promoted tables and figures. The release uses the explicit version tag and MIT License recorded in the package manifest. A permanent public repository URL and DOI are an external deposit step for the corresponding author and are intentionally not fabricated in this manuscript.

\textbf{Funding and correspondence.}\;This work was supported in part by the National Natural Science Foundation of China under Grants 62221005 and 62222601. The authors declare no competing interests. No new human or animal data were collected; the secondary analyses use public, de-identified source records under their stated access conditions. The corresponding author is Shuyin Xia (xiasy@cqupt.edu.cn). Author contributions and related-manuscript disclosures are listed in the accompanying submission metadata.
\begin{thebibliography}{35}\tiny
\bibitem{angelopoulos2023} A. N. Angelopoulos and S. Bates, Conformal prediction: a gentle introduction. \emph{Foundations and Trends in Machine Learning} 16, 494--591 (2023).
\bibitem{athey2016} S. Athey and G. Imbens, Recursive partitioning for heterogeneous causal effects. \emph{PNAS} 113, 7353--7360 (2016).
\bibitem{breiman2001} L. Breiman, Random forests. \emph{Machine Learning} 45, 5--32 (2001).
\bibitem{brier1950} G. W. Brier, Verification of forecasts expressed in terms of probability. \emph{Monthly Weather Review} 78, 1--3 (1950).
\bibitem{clopper1934} C. J. Clopper and E. S. Pearson, The use of confidence or fiducial limits illustrated in the case of the binomial. \emph{Biometrika} 26, 404--413 (1934).
\bibitem{freedman1981} D. Freedman and P. Diaconis, On the histogram as a density estimator. \emph{Z. Wahrscheinlichkeitstheorie} 57, 453--476 (1981).
\bibitem{gneeiting2007} T. Gneiting and A. E. Raftery, Strictly proper scoring rules, prediction, and estimation. \emph{JASA} 102, 359--378 (2007).
\bibitem{guo2017} C. Guo et al., On calibration of modern neural networks. In \emph{ICML}, 1321--1330 (2017).
\bibitem{lei2018} J. Lei et al., Distribution-free predictive inference for regression. \emph{JASA} 113, 1094--1111 (2018).
\bibitem{pedrycz2013} W. Pedrycz, Building the fundamentals of granular computing. \emph{Applied Soft Computing} 13, 4209--4218 (2013).
\bibitem{xial2019} S. Xia et al., Granular ball computing classifiers for efficient, scalable and robust learning. \emph{Information Sciences} 483, 136--152 (2019).
\bibitem{xia2023} S. Xia et al., Granular-ball computing: an efficient, robust, and interpretable adaptive multi-granularity representation and computation method. arXiv:2304.11171 (2023).
\bibitem{xia2025} S. Xia et al., GBRS: A Unified Granular-Ball Learning Model of Pawlak Rough Set and Neighborhood Rough Set. \emph{IEEE TNNLS} 36, 1719--1733 (2025).
\bibitem{angryk2020} R. A. Angryk et al., Multivariate time series dataset for space weather data analytics. \emph{Scientific Data} 7, 227 (2020), doi:10.1038/s41597-020-0548-x; data doi:10.7910/DVN/EBCFKM.
\bibitem{goldberger2000} A. L. Goldberger et al., PhysioBank, PhysioToolkit, and PhysioNet. \emph{Circulation} 101, e215--e220 (2000).
\bibitem{schalk2009eegmmidb} G. Schalk, EEG Motor Movement/Imagery Dataset v1.0.0. \emph{PhysioNet} (2009), doi:10.13026/C28G6P.
\bibitem{pedregosa2011} F. Pedregosa et al., Scikit-learn: Machine learning in Python. \emph{Journal of Machine Learning Research} 12, 2825--2830 (2011).
\bibitem{natarajan2013} N. Natarajan et al., Learning with noisy labels. In \emph{NeurIPS} 26 (2013).
\bibitem{romano2019} Y. Romano, E. Patterson and E. Candes, Conformalized quantile regression. In \emph{NeurIPS} 32 (2019).
\bibitem{chernozhukov2018} V. Chernozhukov et al., Double/debiased machine learning for treatment and structural parameters. \emph{Econometrics Journal} 21, C1--C68 (2018).
\bibitem{mousavi2019stead} M. Mousavi et al., STanford EArthquake Dataset (STEAD): A global data set of seismic signals for AI. \emph{IEEE Access} 7, 179464--179476 (2019), doi:10.1109/ACCESS.2019.2947848.
\bibitem{kasieczka2021} G. Kasieczka et al., The LHC Olympics 2020: A community challenge for anomaly detection in high energy physics. \emph{Reports on Progress in Physics} 84, 124201 (2021), doi:10.1088/1361-6633/ac3c3b.
\bibitem{bump_hunter} L. Anzalone et al., pyBumpHunter: a model-independent bump-hunting tool for high-energy physics. arXiv:2208.14760 (2022).
\bibitem{navon1977} D. Navon, Forest before trees. \emph{Cognitive Psychology} 9, 353--383 (1977).
\bibitem{schyns1994} P. G. Schyns and A. Oliva, From blobs to boundary edges. \emph{Psychological Science} 5, 195--200 (1994).
\end{thebibliography}
\end{document}
